\chapter{Chiral coordinate changes and a compensated quotient}
\label{chap:cc-chiral-collision}

The change of coordinates on a cotangent bundle transforms both positions and momenta.
For differential operators on the projective line, it can also contain a scalar period:
\[
 x'=x^{-1},\qquad v'=-x^2v+\nu x,\qquad [v,x]=1.
\]
Here $x$ acts by multiplication, $v$ acts by differentiation, and $\nu$ is central.
These formulas describe operators at a single energy.
We now ask for fields whose lowest-energy operators obey these relations.
The fields must retain singular products and derivatives with respect to their position.
That additional requirement changes the coordinate formula.

There is a second way to seek these fields.
Start with two coordinate--momentum pairs and divide by simultaneous scaling of the coordinates.
The classical scaling symmetry becomes anomalous for the fields.
A compensating field removes the anomaly, but it also changes the surviving current algebra.
We will compute the quotient, including the surviving field and the ghost zero mode.

Both constructions fit a single family.
Its parameter controls a double pole, whereas $\nu$ controls the lowest-energy action of a field.
The family will have the same differential-operator quotient for every value of the double-pole parameter.
It therefore gives an explicit example of information lost by passing from fields to lowest-energy operators.

\section{Fields made from polynomial states}

Work over $B=\mathbb C[\rho]$, where $\rho$ is an indeterminate.
The letter $z$ denotes the formal position of a field on a complex curve.
It is independent of the target coordinate introduced below.
In particular, two target coordinates will not mean two complex dimensions for $z$.

A state space is a $B$-module $V$ with a distinguished element $1$, called the vacuum.
A field on $V$ is a series
\[
 A(z)=\sum_{n\in\mathbb Z}A_{(n)}z^{-n-1}
\]
of endomorphisms such that $A_{(n)}u=0$ for each fixed $u$ and all sufficiently large $n$.
The operators $A_{(n)}$ are its modes.
All series below are formal.
No norm convergence or completion of the state space is required.

Take the polynomial state space
\[
 V_\rho=B[g_n,b_n,a_n\mid n\geq0].
\]
Give $g_n$ weight $n$ and give $b_n,a_n$ weight $n+1$.
The vacuum is the constant polynomial $1$.
Define translation by the derivation
\[
 T=\sum_{n\geq0}(n+1)
 \left(g_{n+1}\frac{\partial}{\partial g_n}
       +b_{n+1}\frac{\partial}{\partial b_n}
       +a_{n+1}\frac{\partial}{\partial a_n}\right).
\]
Every polynomial uses finitely many variables, so this sum defines an operator.
The three generating fields are
\begin{align}
 \gamma(z)&=\sum_{n\geq0}g_nz^n
             -\sum_{n\geq0}\frac{\partial}{\partial b_n}z^{-n-1},\label{eq:cc-free-fields}\\
 \beta(z)&=\sum_{n\geq0}b_nz^n
             +\sum_{n\geq0}\frac{\partial}{\partial g_n}z^{-n-1},\notag\\
 a(z)&=\sum_{n\geq0}a_nz^n
             +2\rho\sum_{n\geq0}(n+1)
                \frac{\partial}{\partial a_n}z^{-n-2}.\notag
\end{align}
The corresponding states are $g_0,b_0,a_0$.
We also denote these states by $\gamma,\beta,a$.
Translation satisfies $[T,A(z)]=\partial_zA(z)$ for each generating field.
For example, shifting the creation index gives the nonnegative powers in this identity.
Commuting $T$ past $\partial/\partial g_n$ gives
$-n\,\partial/\partial g_{n-1}$ and supplies the negative powers.

Normal ordering moves multiplication operators to the left of differentiation operators.
For two fields it is written $:A(z)B(z):$.
For a monomial in generating fields and their derivatives, full normal ordering moves all its differentiation operators to the right simultaneously.
This convention defines the field of every polynomial state.
For example, the field of $g_0^2b_0$ is the full normal product $:\gamma^2\beta:$.
The field of $g_n$ is $T^n\gamma/n!$, and similarly for $b_n,a_n$.

The singular product $A(z)B(w)\sim\sum_{j\geq0}(A_{(j)}B)(w)(z-w)^{-j-1}$ lists its negative powers of $z-w$.
The symbol $\sim$ discards its regular powers.
The binary normal product $:AB:$ is the regular coefficient $A_{(-1)}B$.
Binary normal products of composite fields need not equal their full normal products.
Nested binary products will always retain their parentheses.

Commuting one differentiation operator past one multiplication operator in
\eqref{eq:cc-free-fields} gives
\begin{equation}\label{eq:cc-basic-poles}
 \beta(z)\gamma(w)\sim\frac1{z-w},\qquad
 \gamma(z)\beta(w)\sim-\frac1{z-w},\qquad
 a(z)a(w)\sim\frac{2\rho}{(z-w)^2}.
\end{equation}
All other generating singular products vanish.
For the first formula, the relevant sum is
$\sum_{n\geq0}w^nz^{-n-1}$, the expansion of $(z-w)^{-1}$ in nonnegative powers of $w/z$.
Differentiating this series gives the last formula.
The pair $(\beta,\gamma)$ is the bosonic beta--gamma system.
The field $a$ generates a Heisenberg algebra with coefficient $2\rho$.
At $\rho=0$ it instead generates a commutative differential algebra: all its singular products vanish.

\begin{lemma}\label{lem:cc-vertex-construction}
The state-field map just defined, together with $1$ and $T$, makes $V_\rho$ a vertex algebra.
The construction also works on $V_\rho[g_0^{-1}]$ and on every specialization $\rho\mapsto r\in\mathbb C$.
\end{lemma}
\begin{proof}
We use the vacuum, translation, and locality definition of a vertex algebra.
The vacuum field is the identity.
Each field applied to $1$ has no negative powers and has the prescribed state as constant coefficient.
The translation identity extends from the generators to their derivatives and full normal products.
Locality means that some power of $z-w$ kills the difference between the two orders of any two fields.

To prove it, commute the differentiation operators in one full normal product past the multiplication operators in the other.
Each commutation either leaves the factors unpaired or replaces one pair by a scalar series from \eqref{eq:cc-basic-poles}.
Repeating gives the sum over all partial pairings of factors from opposite products.
Different orders of these commutations give the same pairing with the same coefficient.
Each pairing contributes a product of derivatives of $(z-w)^{-1}$ or $(z-w)^{-2}$.
There are finitely many pairings for fixed polynomial states.
The two orders are expansions of the same expression, with the sign in the second formula of \eqref{eq:cc-basic-poles}.
Multiplication by a sufficiently high power of $z-w$ removes all denominators and makes the expansions equal.
This proves locality.

For localization, expand the inverse field around its multiplication part:
\[
 \gamma(z)^{-1}=\sum_{j\geq0}(-1)^j
 \left(\sum_{n\geq0}g_nz^n\right)^{-j-1}
 \left(-\sum_{n\geq0}\frac{\partial}{\partial b_n}z^{-n-1}\right)^j.
\]
On a fixed state, only finitely many values of $j$ act nontrivially.
The first factor has a power-series inverse because $g_0$ is invertible.
The same prescription defines rational functions of $\gamma$ whose denominators are powers of $\gamma$.
For two localized states, every contraction still uses a beta or Heisenberg factor.
Their total number bounds the number of contractions, even when differentiating a Laurent function of $\gamma$.
The preceding locality argument therefore still applies.
No step divides by $\rho$.

For later use, the products satisfy the iterate identity
\begin{align}
 (A_{(m)}B)_{(n)}={}&\sum_{i\geq0}(-1)^i\binom mi
 \bigl(A_{(m-i)}B_{(n+i)}
       -(-1)^m B_{(m+n-i)}A_{(i)}\bigr).
 \label{eq:cc-iterate}
\end{align}
For an arbitrary integer $m$, the binomial coefficient means
$\binom mi=m(m-1)\cdots(m-i+1)/i!$, with $\binom m0=1$.
Indeed, multiply the difference of the two field orders by $(z-w)^m$ and take the coefficient defining $A_{(m)}B$.
Expand $(z-w)^m$ in each order by the binomial formula.
Taking the remaining $w$ coefficient gives exactly the two sums in \eqref{eq:cc-iterate}.
On any fixed state the required coefficients are finite.
Applying this argument to three fields is legitimate because their pairings are scalar.
Each partial pairing occurs in both iterated expansions with the same coefficient.
Thus the identity holds for the full state-field map, including localized states.
\end{proof}

The multiplication and differentiation realization supplies more than a table of poles.
It specifies every composite state, every regular coefficient, and all iterated products.
These data will determine whether a coordinate change has an inverse.

\section{A family of coordinate changes}

Set $\kappa=\rho-2$ and define three weight-one fields
\begin{equation}\label{eq:cc-currents}
 e=\beta,\qquad h=-2:\gamma\beta:+a,\qquad
 f=-:\gamma^2\beta:+:a\gamma:+\kappa T\gamma.
\end{equation}
The derivative term has the same weight as the other terms in $f$.
It is forced by the singular products.
For example, omitting it at $\rho=0$ gives a field $f_0$ with
\[
 f_0(z)f_0(w)\sim
 -\frac{4\gamma(w)^2}{(z-w)^2}
 -\frac{4\gamma(w)T\gamma(w)}{z-w}.
\]
A new momentum must have zero singular product with itself.
The ordinary cotangent formula therefore cannot serve as the field transformation.

\begin{proposition}\label{prop:cc-current-poles}
The fields in \eqref{eq:cc-currents} have singular products
\begin{align*}
 h(z)e(w)&\sim\frac{2e(w)}{z-w},&
 h(z)f(w)&\sim-\frac{2f(w)}{z-w},\\
 e(z)f(w)&\sim\frac{\kappa}{(z-w)^2}+\frac{h(w)}{z-w},&
 h(z)h(w)&\sim\frac{2\kappa}{(z-w)^2},\\
 e(z)e(w)&\sim0,& f(z)f(w)&\sim0.
\end{align*}
These are the affine $\mathfrak{sl}_2$ current relations at level $\kappa$.
Here $\mathfrak{sl}_2$ means the three-dimensional Lie algebra with
$[h,e]=2e$, $[h,f]=-2f$, and $[e,f]=h$.
The level is the coefficient of the double pole in $e(z)f(w)$.
\end{proposition}
\begin{proof}
Put $r=z-w$.
In $e(z)f(w)$, contraction with $-\gamma^2\beta$ gives $-2\gamma\beta/r$.
Contraction with $a\gamma$ gives $a/r$.
Contraction with $\kappa T\gamma$ gives $\kappa/r^2$.
This proves the third formula.
In $h(z)h(w)$, the two beta--gamma contractions give $-4/r^2$.
The Heisenberg contraction adds $2\rho/r^2$.

The remaining formulas follow from the same pairing rule with at most two contractions.
The potentially nonzero self-product can be seen explicitly before adding the derivative.
For $f_0=-\gamma^2\beta+a\gamma$, its double and simple poles are
\[
 \frac{2(\rho-2)\gamma^2}{r^2}
 +\frac{2(\rho-2)\gamma T\gamma}{r}.
\]
The sum of the two cross-products with $\kappa T\gamma$ is their negative.
The derivative has zero self-product.
Thus $ff$ has no poles.
For $hf$, the double-pole contributions from the first two summands in $f$ cancel those from its derivative term.
Its simple pole is
$2\gamma^2\beta-2a\gamma-2\kappa T\gamma=-2f$.
Finally, $he$ has its single beta--gamma contraction, and $ee$ has none.
\end{proof}

On the overlap where $\gamma$ is invertible, put
\begin{equation}\label{eq:cc-transition}
 \gamma'=\gamma^{-1},\qquad
 \beta'=f,\qquad
 a'=a+2\rho\gamma^{-1}T\gamma.
\end{equation}
The logarithmic derivative $\gamma^{-1}T\gamma$ is an element of the localized state algebra.
No logarithm of $\gamma$ has been adjoined.

\begin{theorem}\label{thm:cc-gluing}
The formulas \eqref{eq:cc-transition} give an isomorphism of localized vertex algebras.
Its inverse is the same formula in the primed fields.
They glue two affine charts to a sheaf $\mathcal V_\rho$ of vertex algebras on $\mathbb P^1$.
The fields $e,h,f$ extend globally, with transition $(e,h,f)\mapsto(f,-h,e)$.
\end{theorem}
\begin{proof}
The pairing rule gives
\[
 \beta'(z)\gamma'(w)\sim r^{-1},\qquad
 a'(z)a'(w)\sim2\rho r^{-2},\qquad
 a'(z)\beta'(w)\sim a'(z)\gamma'(w)\sim0.
\]
The self-products of $\beta'$ and $\gamma'$ vanish.
For the mixed Heisenberg--momentum product, write $t=\gamma^{-1}T\gamma$.
The contraction $a(z)f(w)$ is $2\rho\gamma(w)/r^2$.
The product $t(z)f(w)$ is $-\gamma(w)/r^2$, with zero simple pole.
These terms cancel in $a'f$.
All primed generating products therefore agree with \eqref{eq:cc-basic-poles}.
The field construction extends the assignment to a homomorphism: substitute the new fields into each full normal product and derivative.
The pairing and residue identities ensure preservation of every product.

To check its inverse, normal ordering must be retained.
For any Laurent function $u(\gamma)$ and polynomial or Laurent function $P(\gamma)$, the regular coefficient is
\begin{equation}\label{eq:cc-regular-correction}
 :u(\gamma)(P(\gamma)\beta):
 =uP\beta-u''P\,T\gamma.
\end{equation}
Indeed, the only cross-contraction is
$-u'(\gamma(z))P(\gamma(w))/r$.
Its first Taylor coefficient is the second term in \eqref{eq:cc-regular-correction}.
Applying this twice gives
\begin{align*}
 :\gamma'\beta':&=-\gamma\beta+a+\rho\gamma^{-1}T\gamma,\\
 :\gamma'(:\gamma'\beta':):&=
 -\beta+\gamma^{-1}a+(\rho+2)\gamma^{-2}T\gamma.
\end{align*}
Consequently
\[
 -:\gamma'(:\gamma'\beta':):+:a'\gamma':
       +(\rho-2)T\gamma'=\beta.
\]
Also $a'+2\rho(\gamma')^{-1}T\gamma'=a$ and $(\gamma')^{-1}=\gamma$.
Thus the homomorphism is involutive.
The first regular-product formula gives $h'=-h$.
The inverse formula gives $f'=e$, while $e'=f$ by definition.

On either affine chart, further principal open subsets are defined by localizing its coordinate algebra.
The same construction localizes the fields.
Sections on a union are pairs of sections agreeing under the overlap isomorphism.
The inverse calculation proves the cocycle condition for this two-chart cover.
This constructs the asserted sheaf and its global currents.
\end{proof}

At $\rho=0$, the field $a$ is central and is unchanged between charts.
At $\rho\ne0$, it has a nonzero double pole and changes by a logarithmic derivative.
The correction in $\beta'$ and the correction in $a'$ are parts of the same coordinate change.

\section{The quadratic current and the energy operator}

The singular products of $e,h,f$ suggest the quadratic expression
\[
 S=\tfrac12:hh:+:ef:+:fe:.
\]
Every colon here denotes a binary regular coefficient.
The expression would lose a derivative term if its factors were multiplied as commuting symbols.

\begin{proposition}\label{prop:cc-sugawara}
The exact identity and singular self-product are
\begin{align}
 S&=2\rho:\beta T\gamma:+\tfrac12:a^2:-Ta,\label{eq:cc-miura}\\
 S(z)S(w)&\sim
 \frac{6\rho(\rho-2)}{(z-w)^4}
 +\frac{4\rho S(w)}{(z-w)^2}
 +\frac{2\rho TS(w)}{z-w}.\notag
\end{align}
At $\rho=0$, $S$ has zero singular product with every local field.
If $\rho$ is invertible, $L=S/(2\rho)$ is a stress tensor of central charge $3-6/\rho$.
\end{proposition}
\begin{proof}
Full normal products on the right give the three regular coefficients
\begin{align*}
 :hh:&=4\gamma^2\beta^2-4a\gamma\beta+a^2
             +4\beta T\gamma-4\gamma T\beta,\\
 :ef:&=-\gamma^2\beta^2+a\gamma\beta+(\rho-2)\beta T\gamma,\\
 :fe:&=-\gamma^2\beta^2+a\gamma\beta+\rho\beta T\gamma
             +2\gamma T\beta-Ta.
\end{align*}
For example, the single contraction in $f(z)e(w)$ has numerator
$2\gamma(z)\beta(z)-a(z)$.
Its first Taylor coefficient contributes $2\beta T\gamma+2\gamma T\beta-Ta$.
This explains the difference between the last two lines.
Their sum with half the first line proves \eqref{eq:cc-miura}.

Put $U=:\beta T\gamma:$ and $A=\tfrac12:a^2:-Ta$.
These fields commute with one another.
Direct contractions give
\begin{align*}
 U(z)U(w)&\sim r^{-4}+2U(w)r^{-2}+TU(w)r^{-1},\\
 A(z)A(w)&\sim(2\rho^2-12\rho)r^{-4}
                    +4\rho A(w)r^{-2}+2\rho TA(w)r^{-1}.
\end{align*}
In the second line the fourth pole from $(-Ta)(z)(-Ta)(w)$ is $-12\rho r^{-4}$.
This derivative contraction accounts for the shift in the central charge.
Since $S=2\rho U+A$, the displayed self-product follows.

A stress tensor is a field $L$ whose self-product is
$c/(2r^4)+2L/r^2+TL/r$ and whose simple-pole action is translation.
Here the generating products are
\[
 L(z)\gamma(w)\sim\frac{T\gamma}{r},\qquad
 L(z)\beta(w)\sim\frac{\beta}{r^2}+\frac{T\beta}{r},\qquad
 L(z)a(w)\sim\frac2{r^3}+\frac a{r^2}+\frac{Ta}{r}.
\]
The simple-pole action extends to all products as a derivation, so it equals $T$ on every state.
The double-pole action gives the weights already assigned.
A field of weight $d$ is primary when its product with $L$ has only the poles
$dA/r^2+TA/r$.
The third pole therefore means that $a$ is not primary for this stress tensor.
At $\rho=0$, formula \eqref{eq:cc-miura} involves only the central field $a$ and its derivatives.
Division by $\rho$ is then unavailable, but centrality of $S$ remains an exact statement.
\end{proof}

Thus the critical level $\kappa=-2$ is a singular point of the conformal construction.
The quadratic expression remains defined there.
Its role changes from generating the stress tensor to generating central operations.

\section{Lowest-energy operators and the Zhu quotient}

We next recover the differential-operator chart from the full field algebra.
Keeping the regular product alone would not give an associative multiplication.
The appropriate operation includes the singular coefficients selected by energy.

For homogeneous states $A,B$, let $d_A$ denote the nonnegative integer weight of $A$.
Define
\[
 A*B=\sum_{j=0}^{d_A}\binom{d_A}{j}A_{(j-1)}B,\qquad
 A\circ B=\sum_{j=0}^{d_A}\binom{d_A}{j}A_{(j-2)}B.
\]
Let $O(V)$ be the linear span of the states $A\circ B$.
We will prove directly for the present free fields that $*$ induces an associative product on $V/O(V)$.
This associative algebra is the Zhu quotient.
The definition uses the weight grading even when $\rho=0$ and the preceding stress tensor is unavailable.
A module assigns fields to the same states on another vector space, preserving the vacuum field and the product residue identities.
An energy grading on that space assigns an integer to each homogeneous vector.
The energy change of a mode is the difference between its output and input energies.

\begin{lemma}\label{lem:cc-zero-mode}
Let a module have nonnegative integral energies, with $A_{(n)}$ changing energy by $d_A-n-1$.
On its energy-zero space, put $o(A)=A_{(d_A-1)}$.
Then $o(A*B)=o(A)o(B)$ and $o(A\circ B)=0$.
\end{lemma}
\begin{proof}
Write $a=d_A$, $b=d_B$.
For the $j$th summand of $A*B$, apply \eqref{eq:cc-iterate} with
$m=j-1$, $n=a+b-j-1$.
The first term contains $B_{(a+b-j-1+i)}$ on the right.
On energy zero it vanishes unless $j=a$ and $i=0$.
That remaining term is $o(A)o(B)$.
In the second term, the index of $B$ is $a+b-2-i$, independent of $j$.
Also $A_{(i)}$ vanishes on energy zero for $i\geq a$.
For $0\leq i<a$, its summed coefficient is, up to a common sign,
\[
 \sum_{j=0}^a(-1)^j\binom aj\binom{j-1}{i}=0.
\]
The equality is the $a$th finite difference of a polynomial of degree $i<a$.
It follows by applying $(1-u)^a$ to the generating polynomials for finite differences, or by induction on $a$.

For $A\circ B$, use $m=j-2$, $n=a+b-j$.
Every first term now annihilates energy zero.
The second terms have the same vanishing finite difference with $\binom{j-2}{i}$.
This proves both assertions, including $a=0$, when the second sums act by zero without any finite difference.
\end{proof}

\begin{theorem}\label{thm:cc-zhu}
On either affine chart,
\[
 V_\rho/O(V_\rho)\cong
 B[\nu]\langle x,v\rangle/(vx-xv-1),
 \qquad [\gamma]\mapsto x,\quad[\beta]\mapsto v,\quad[a]\mapsto\nu.
\]
The localized quotient permits $x^{-1}$.
These quotients glue under
\[
 x'=x^{-1},\qquad v'=-x^2v+\nu x,\qquad \nu'=\nu.
\]
In particular, the sheafified Zhu quotient of $\mathcal V_\rho$ is independent of $\rho$.
\end{theorem}
\begin{proof}
Construct a module with energy-zero space $M_0=B[\nu,x]$.
For the coordinate and momentum zero modes use $x$ and $\partial_x$.
For the Heisenberg zero mode use multiplication by $\nu$.
Adjoin freely all creation modes of positive energy.
The positive modes act as the corresponding differentiations, with commutators prescribed by \eqref{eq:cc-basic-poles}.
Writing $\mathsf a_m=a_{(m)}$ for the Heisenberg modes, their commutator is
$[\mathsf a_m,\mathsf a_n]=2\rho m\delta_{m+n,0}$.
Its zero mode therefore commutes with all modes, as required.
The same scalar-pairing construction proves the module identities.
The module is graded in nonnegative energies, with precisely $M_0$ in energy zero.

The preceding lemma defines a linear map $\pi:V\to\operatorname{End}_B(M_0)$ which kills $O(V)$ and preserves $*$.
For a full normal state $f(\gamma)\beta^p a^q$, its value is
\begin{equation}\label{eq:cc-top-action}
 \pi\bigl(f(\gamma)\beta^p a^q\bigr)
       =f(x)\partial_x^p\nu^q.
\end{equation}
Indeed, a normal product puts every positive mode to the right, where it kills $M_0$.
A surviving negative mode would raise energy.
Only individual zero modes can therefore contribute to the total zero mode on $M_0$.
Their order is the displayed multiplication followed by differentiation.

It remains to show that these states span the quotient.
Filter states by the number of beta and Heisenberg factors, counting their derivatives with the same multiplicity.
All contractions lower this filtration.
Consequently the highest-filtration part of $A\circ B$ is
\[
 (TA+d_AA)B
\]
in the commutative polynomial algebra of jets.
For $A=T^n\gamma,T^n\beta,T^na$, these relations reduce every higher jet to
\[
 T^n\gamma=0\ (n\geq1),\qquad
 T^n\beta=(-1)^n n!\beta,\qquad
 T^na=(-1)^n n!a
\]
in that highest-filtration quotient.
These are triangular relations with leading coefficient one.
Multiplying them by arbitrary jet monomials reduces every state to a combination of the states in \eqref{eq:cc-top-action}, plus lower-filtration terms.
Repeat at the lower filtration.
At filtration zero there are no contractions, so the reduction terminates exactly.
This argument uses only elements $A\circ B$ and their highest parts.
It does not assume in advance that $O(V)$ is an ideal.

The operators in \eqref{eq:cc-top-action} are linearly independent over $B$.
For a relation with largest differential order $N$, commute it $N$ times with multiplication by $x$.
The result is $N!$ times its leading coefficient, which must vanish.
Descending in $N$ proves all coefficients zero, including their distinct powers of $\nu$.
Thus $\ker\pi=O(V)$.
Since $\pi$ preserves $*$, this also proves the ideal property and associativity of the quotient.
The commutator is $[v,x]=1$ by its action on polynomials.

The same argument works over $B[\nu,x,x^{-1}]$.
For every homogeneous state, $A\circ1=TA+d_AA$.
Thus $[T\gamma]=0$ and $[Ta]=-[a]$.
The filtration argument also gives $[\gamma^{-1}T\gamma]=0$.
Applying \eqref{eq:cc-top-action} to the full normal terms in \eqref{eq:cc-transition} gives the asserted overlap map.
Define the differential-operator sheaf by these two Weyl charts and this overlap.
The local isomorphisms agree on all further localizations, and hence identify the sheafified quotients.
No assertion about commuting the Zhu functor with global sections has been used.
\end{proof}

The images of the global currents are
\[
 E=v,\qquad H=\nu-2xv,\qquad F=-x^2v+\nu x.
\]
Their commutators are $[H,E]=2E$, $[H,F]=-2F$, and $[E,F]=H$.
The central quadratic expression is
\[
 \Omega=H^2+2H+4FE=\nu(\nu+2).
\]
For completeness, use $vx=xv+1$ to expand
$H^2=\nu^2-4\nu xv+4x^2v^2+4xv$ and
$4FE=-4x^2v^2+4\nu xv$.
Adding $2H=2\nu-4xv$ gives the displayed scalar.

The full field calculation gives the same answer:
\[
 2[S]=[a^2]-2[Ta]=\nu^2+2\nu.
\]
Here $[\beta T\gamma]=0$ by \eqref{eq:cc-top-action} and the preceding jet reduction.
Moreover, $e*f=:ef:+h$ and $f*e=:fe:-h$.
Therefore the Zhu image of $2S$ is precisely $H^2+2(EF+FE)=\Omega$.
This proves compatibility of the quadratic operation with the operator quotient.
It also exhibits the loss of information: the same scalar $\nu(\nu+2)$ comes from every level $\kappa=\rho-2$.

\section{The scaling anomaly and its compensation}

We now construct a differential which implements the scaling quotient of two coordinate--momentum pairs.
Take even fields $(g,p)$ and $(y,q)$, with
\[
 p(z)g(w)\sim r^{-1},\qquad q(z)y(w)\sim r^{-1}.
\]
Their weights are $0,1$ and $0,1$, respectively.
Their state space and localized state spaces are defined by the same polynomial mode construction as before.
The infinitesimal scaling current is
\[
 J=:gp:+:yq:.
\]
Its simple-pole action multiplies $g,y$ by $1$ and $p,q$ by $-1$.
Double contraction gives $J(z)J(w)\sim-2r^{-2}$.

Introduce odd fields $b,c$ of weights $1,0$ with
$b(z)c(w)\sim c(z)b(w)\sim r^{-1}$.
Their state space is an exterior algebra of creation modes.
Odd differentiation replaces ordinary differentiation in the construction of their fields.
Interchanging two odd factors introduces a minus sign.
These auxiliary fields are the antighost and ghost.

For a weight-one odd field $j_Q$, its zero mode is the odd operator
$\operatorname{Res}_zj_Q(z)$, where the residue selects the coefficient of $z^{-1}$.
The iterate identity shows that a zero mode is a derivation of every field product.
The proposed charge $Q_0=\operatorname{Res}:cJ:$ has
\[
 Q_0b=J,\qquad Q_0J=-2Tc,\qquad Q_0^2b=-2Tc\ne0.
\]
Thus the anomaly prevents this charge from defining a cochain complex.
The derivative $Tc$ is a nonzero creation state.

Adjoin an independent even field $j$ with $j(z)j(w)\sim2r^{-2}$.
Denote the resulting state space by $\mathsf C$.
It is the tensor product of the two polynomial pair spaces, the polynomial Heisenberg space, and the exterior ghost space.
Set
\begin{equation}\label{eq:cc-brst-charge}
 K=J+j,\qquad Q=\operatorname{Res}:cK:.
\end{equation}
The field $K$ has zero singular self-product.
Also $c$ has zero self-product and commutes with $K$.
It follows that $:cK:$ has zero singular self-product, so $Q^2=0$.
This is the BRST differential, an odd differential defined by the charge in \eqref{eq:cc-brst-charge}.
Give $c$ cochain degree $1$, $b$ degree $-1$, and the even fields degree zero.
The differential raises cochain degree by one and preserves energy.
Its generating actions are
\begin{align*}
 Qg&=:cg:,& Qy&=:cy:,& Qp&=-:cp:,& Qq&=-:cq:,\\
 Qj&=2Tc,& Qb&=K,& Qc&=0,& QK&=0.
\end{align*}

\section{The complete local quotient}

Let $\mathsf C_g=\mathsf C[g_0^{-1}]$ be the localization at the coordinate state of $g$.
Define
\begin{equation}\label{eq:cc-residual-fields}
 x=y/g,\qquad v=:gq:,\qquad t=g^{-1}Tg,\qquad A=-j+2t.
\end{equation}
The first two fields are the coordinate and momentum of the projective chart.
The third is a logarithmic derivative in the scaling direction.
It satisfies $Qt=Tc$, so $A$ is closed.
The singular products are
\begin{align}
 v(z)x(w)&\sim r^{-1},& A(z)A(w)&\sim2r^{-2},\label{eq:cc-residual-poles}\\
 A(z)x(w)&\sim0,& A(z)v(w)&\sim0.\notag
\end{align}
The field $K$ has zero singular product with $x,v,A$, whereas
$K(z)t(w)\sim r^{-2}$ and $K(z)g(w)\sim g(w)r^{-1}$.
In particular, the extra field $A$ is not central.

Closed generators give classes in cohomology.
To prove that they generate all cohomology, we must also remove every state in the complementary directions by an explicit homotopy.

For a field of weight $d$, use physical indices $u_m=u_{(m+d-1)}$.
Then $u_m$ changes energy by $-m$.
The modes $K_{-r},t_{-r},A_{-r},v_{-r}$ for $r\geq1$ are creation operators.
For the weight-zero fields, the creation operators include $x_0,c_0$ and the modes $x_{-r},c_{-r}$ for $r\geq1$.

\begin{theorem}\label{thm:cc-brst-cohomology}
The cohomology of the localized complex is
\[
 H^\bullet(\mathsf C_g,Q)
 \cong V_{\beta\gamma}(x,v)\otimes\mathcal H_2(A)\otimes\Lambda(c_0).
\]
Here $V_{\beta\gamma}(x,v)$ is the polynomial beta--gamma vertex algebra,
$\mathcal H_2(A)$ is the Heisenberg vertex algebra with double-pole coefficient $2$,
and $\Lambda(c_0)$ is the exterior algebra on one odd degree-one class.
Translation acts by zero on $[c_0]$.
On the relative subcomplex $\ker b_0\cap\ker K_0$, the cohomology is
$V_{\beta\gamma}(x,v)\otimes\mathcal H_2(A)$.
\end{theorem}
\begin{proof}
First establish an ordered basis of states.
The variables $g_0^{\pm1}$ and the jets of $t$ replace all positive jets of $g$.
Indeed, $Tg=tg$, and its successive derivatives express $T^ng$ triangularly in these new variables.
The leading coefficient in the reverse expression for $T^{n-1}t$ is $g_0^{-1}$.
Similarly, $y=xg$ replaces every jet of $y$ by the jets of $x$ and $g$.

Filter the original states by the number of factors of $p,q,j,b$, including their derivatives.
Put the coordinate jets and $c$ in filtration degree zero.
Every contraction lowers this filtration.
At its highest part, the change from momenta to $v,A,K$ is
\[
 v=gq,\qquad A=-j,\qquad K=gp+yq+j.
\]
Its determinant has absolute value $g^2$, which is invertible after localization.
Differentiating this change gives triangular changes at every jet order, with the same invertible diagonal blocks.
Thus the associated graded algebra has a polynomial basis in the new variables, with exterior ghost variables.

Fix an order for their creation modes and put $g_0^n$, $n\in\mathbb Z$, last.
The corresponding ordered states form a basis of the original filtered vector space.
Spanning follows by replacing the highest-filtration symbol and descending in filtration.
Independence follows by taking the highest nonzero symbol of a finite relation.
Each descent is finite because each state has finite filtration degree.
This establishes a vector-space identification, with residual states on the left, with
\begin{equation}\label{eq:cc-brst-basis}
 V_{\beta\gamma}(x,v)\otimes\mathcal H_2(A)
 \otimes\mathbb C[K_{-r},t_{-r}\mid r\geq1]
 \otimes\Lambda(b_{-r},c_{-r}\mid r\geq1)
 \otimes\mathbb C[g_0,g_0^{-1}]\otimes\Lambda(c_0).
\end{equation}
It does not assert that all these factors commute as vertex algebras.
In particular, $K$ acts nontrivially on $g$.
The stated ordering retains that action.

Since $Q$ is a zero-mode derivation, its action on each ordered state is determined by its commutators with these modes.
The field equations give
\[
 [Q,b_{-r}]_+=K_{-r},\qquad
 [Q,t_{-r}]=r c_{-r},\qquad Q(g_0^n)=n c_0g_0^n.
\]
It kills the residual fields, $K$, and $c$.
The symbol $[\ ,\ ]_+$ denotes the anticommutator of two odd operators.
Moving $c_0$ into the fixed order introduces only exterior signs.
It has zero anticommutator with $b_{-r}$ for $r\geq1$.
All other replacements keep the fixed creation order.
Therefore, on the vector space \eqref{eq:cc-brst-basis}, the differential is exactly
\begin{equation}\label{eq:cc-brst-polynomial-differential}
 d=\sum_{r\geq1}\left(
 K_{-r}\frac{\partial}{\partial b_{-r}}
 +r c_{-r}\frac{\partial}{\partial t_{-r}}\right)
 +c_0g_0\frac{\partial}{\partial g_0}.
\end{equation}
Odd derivatives are left exterior derivatives.

Let $N$ count the total polynomial and exterior degree in the four paired variables for all $r\geq1$.
Define the odd operator
\[
 h=\sum_{r\geq1}\left(
 b_{-r}\frac{\partial}{\partial K_{-r}}
 +\frac1r t_{-r}\frac{\partial}{\partial c_{-r}}\right).
\]
The product rule gives $dh+hd=N$.
Terms with different variables anticommute, and the $c_0$ term anticommutes with $h$.
Thus $h/N$ contracts every positive-$N$ summand.
Only finitely many summands act on any state.

At $N=0$, the remaining differential on the Laurent monomial $g_0^n$ is multiplication by $n c_0$.
For $n\ne0$, the homotopy $n^{-1}\partial/\partial c_0$ contracts this two-term complex.
For $n=0$, the differential vanishes and leaves $1,c_0$.
This proves the claimed vector-space cohomology and its cochain degrees.
The closed residual generators and $c_0$ represent every surviving basis vector.
Their products in cohomology are \eqref{eq:cc-residual-poles}, with $[c_0]$ central and square zero.
Also $Tc=Qt$, so translation kills $[c_0]$.
They therefore define a homomorphism from the asserted tensor product into cohomology.
The vector-space calculation proves that this homomorphism is an isomorphism.
There is no translation-preserving inclusion of the constant exterior factor into the original complex: $Tc$ vanishes only after passing to cohomology.

Finally, $[Q,b_0]_+=K_0$ and $[Q,K_0]=0$.
Hence $\ker b_0\cap\ker K_0$ is a subcomplex.
In \eqref{eq:cc-brst-basis}, $K_0$ acts by $n$ on $g_0^n$, and $b_0$ differentiates with respect to $c_0$.
The relative subcomplex consequently has $n=0$ and no $c_0$ factor.
The first contraction still applies and leaves exactly the two residual vertex factors.
\end{proof}

The surviving odd class records the scaling direction in absolute cohomology.
It would change if a logarithm of $g_0$ were adjoined, since then $Q\log g_0=c_0$.
The Laurent state space used here contains no such logarithm.
The difference between the absolute and relative quotients is therefore a difference of specified complexes.

\section{Descent of the quotient and its stress tensor}

On the second chart, define
\[
 x'=g/y,\qquad v'=:yp:,\qquad A'=-j+2y^{-1}Ty.
\]
The fields $x',v',A'$ are closed.
On their common domain, the classes obey
\begin{equation}\label{eq:cc-gauged-overlap}
 [x']=[x^{-1}],\qquad
 [v']=[-:x(:xv:):+:Ax:-Tx],\qquad
 [A']=[A+2x^{-1}Tx].
\end{equation}
These are \eqref{eq:cc-transition} at $\rho=1$.
Thus the compensated quotient has affine level $\kappa=-1$.
The critical construction instead has $\rho=0$ and level $-2$.

\begin{proposition}\label{prop:cc-quotient-descent}
The relative local BRST cohomologies glue to $\mathcal V_1$.
Their stress tensor has central charge $-3$ and is induced from the compensated theory.
\end{proposition}
\begin{proof}
The regular-product corrections can be checked while keeping the original fields.
Define a radial momentum
\[
 P=p+(y/g)q+\frac{Tg}{2g^2}.
\]
Contractions give $P(z)g(w)\sim r^{-1}$, $P(z)P(w)\sim0$, and zero singular products of $P$ with $x,v$.
The regular coefficient also gives
\[
 :gP:-\tfrac12t=J.
\]
The half-derivative term is needed for both assertions.
Substituting $y=xg$ and retaining binary products now yields
\begin{align*}
 :yp:-:xK:&=-:x(:xv:):+:Ax:-Tx,\\
 :gp:-:yq:-K&=-2:xv:+A,\\
 -j+2y^{-1}Ty&=A+2x^{-1}Tx.
\end{align*}
These equalities can also be derived directly from
\eqref{eq:cc-regular-correction}, applied to $g^{-1}$ and $y/g$.
Since $:xK:=Q(:xb:)$ and $K=Qb$, they prove \eqref{eq:cc-gauged-overlap} and the corresponding global Cartan field.
The local cohomology theorem supplies all states, not only the displayed generators.
Its isomorphisms localize, so the overlap calculation proves the sheaf assertion.

The original matter stress tensor is
$L_{\mathrm{mat}}=:pTg:+:qTy:$.
The compensated stress tensor, including ghosts, is
\[
 L_{\mathrm{tot}}=L_{\mathrm{mat}}+\tfrac14:j^2:
                     +\tfrac12Tj-:bTc:.
\]
The derivative of $j$ fixes its background charge.
Two beta--gamma pairs contribute central charge $4$.
For $j(z)j(w)\sim2r^{-2}$, the field $\tfrac14:j^2:+\tfrac12Tj$ contributes $1-6=-5$.
The odd pair contributes $-2$.
These numbers follow from the respective double contractions in the fourth pole, which is half the central charge.
The total central charge is therefore $-3$.

On the residual fields put
\[
 L_{\mathrm{res}}=:vTx:+\tfrac14:A^2:-\tfrac12TA.
\]
It is the stress tensor $S/2$ of \eqref{eq:cc-miura} at $\rho=1$.
It is closed because every residual field is closed.
The equality of the stress tensors in cohomology has an explicit primitive:
\begin{equation}\label{eq:cc-stress-primitive}
 L_{\mathrm{tot}}-L_{\mathrm{res}}=Q(:bt:).
\end{equation}
To verify it, the coordinate calculation gives
\[
 L_{\mathrm{mat}}=:PTg:+:vTx:-\tfrac12Tt.
\]
Use $A=-j+2t$ and $:gP:=J+t/2$ in the difference of the two stress tensors.
Its even part becomes $:Kt:$.
Its ghost part is $-:bTc:$.
The odd derivation rule gives
$Q(:bt:)=:Kt:-:bTc:$, proving \eqref{eq:cc-stress-primitive}.
The primitive lies in the relative subcomplex: $b_0$ and $K_0$ both kill it.
Thus the conformal structure descends to the stated relative quotient.
The global current expression for $S$ also proves compatibility of this stress tensor between the two charts.
\end{proof}

There are now two explicit constructions and an exact comparison of what they retain.
The critical sheaf $\mathcal V_0$ has a central weight-one field and a central quadratic field.
The compensated quotient $\mathcal V_1$ has a Heisenberg field and a stress tensor of central charge $-3$.
Their sheafified Zhu quotients have the same Weyl charts, period $\nu$, and quadratic value $\nu(\nu+2)$.
Thus agreement of the operator quotient does not identify the chiral theories.
The differing double poles and the surviving Heisenberg field specify the additional data that such an identification would have to preserve.
